% ============================================================
%  Rasmus Nielsen, Philosophiske Grundproblemer:
%  Festskrift i Anledning af Universitets Firehundredaarsfest
%  Juni 1879.
%  Kjøbenhavn: Gyldendalske Boghandel (Trykt hos J.N. Schultz), 1879.
%  77 pp.
%
%  TRANSCRIPTION (Danish)
%
%  Structure (VERIFIED against the color scan; offset PDF = printed + 9):
%    [Introduction, untitled] (pp. 3--4; p.1 title, p.2 blank)
%    I.   Erkjendelsesproblemet            (pp. 4--26)   heading at foot of p.4
%    II.  Realitetsproblemet               (pp. 27--50)
%    III. De ueensartede Principers Problem (pp. 50--77)  heading on p.50
%  No Indhold/table-of-contents page exists in the original;
%  the \tableofcontents below is an editorial aid.
%  Only three numbered sections. The soul--body problem, the
%  problem of freedom, and the Tro-og-Viden discussion all fall
%  WITHIN section III ("De ueensartede Principers Problem"), which
%  is not separately titled. (The earlier scaffold guesses
%  "III. Frihed, Sjæl og Legeme" and a fourth "Tro og Viden"
%  section were wrong.)
%
%  Key passage: p. 72 — "Principerne for Tro og Viden ere
%  absolut ueensartede" — the most explicit popular statement
%  of the central thesis, concluded with the claim that the
%  dualistic separation is itself the condition for a reconciling
%  unity of faith and knowledge.
% ============================================================

\documentclass[12pt]{article}
\usepackage[utf8]{inputenc}
\usepackage[T1]{fontenc}
\usepackage{libertinus}
\usepackage{libertinust1math}
\usepackage{textalpha}   % polytonic Greek typed directly (e.g. ὕλη ἄμορφος, p.13)
\usepackage[danish]{babel}
\usepackage[protrusion=true,expansion=true]{microtype}
\usepackage[margin=1.4in]{geometry}
\usepackage{setspace}
\usepackage{fancyhdr}
\usepackage[colorlinks=true,linkcolor=black,urlcolor=black,
  pdftitle={Philosophiske Grundproblemer},
  pdfauthor={Rasmus Nielsen}]{hyperref}
\setlength{\headheight}{15pt}
\pagestyle{fancy}
\fancyhf{}
\fancyhead[LE,RO]{\thepage}
\fancyhead[RE]{\textit{Philosophiske Grundproblemer}}
\fancyhead[LO]{\textit{\rightmark}}
\setlength{\parindent}{1.5em}
\setlength{\parskip}{0pt}
\onehalfspacing

\title{\textbf{Philosophiske Grundproblemer}\\[0.5em]
  \large Festskrift i Anledning af Universitets Firehundredaarsfest}
\author{R.\ Nielsen}
\date{Kjøbenhavn: Gyldendalske Boghandel, 1879}

\begin{document}
\maketitle
\tableofcontents
\bigskip


% ============================================================
%  INDLEDNING
% ============================================================

% pp. 3--4  (title on p.1, blank p.2; the untitled introduction begins p.3)

\noindent
Philosophien har ikke som andre Videnskaber den Fordeel at
kunne bygge paa umiddelbart faste Forudsætninger, uomtvistelige
Grundsætninger; thi der gives i den hele Philosophie ikke en
eneste Sætning, der staaer saa fast, at den jo, naar den behandles
for sig, kan angribes af Kritik og opløses i Dialektik. Af den
Grund kan Philosophien da heller ikke ret vel udstykke sine
Beviser, saaledes som f.\,Ex.\ Mathematiken. Endskjøndt de
matematiske Beviser maae foredrages i stræng Orden, naar
fuld, videnskabelig Indsigt skal opnaaes, saa udsætter
Mathematikeren sig dog ikke for nogen væsenlig Misforstaaelse,
hvis han i et eller andet Øiemed behandler en særlig Opgave
isoleret. De Kyndige vide, hvorledes Forudsætningerne kunne
bevises, og de, hvis Forkundskaber ikke strække til, lade sig
nøie med Beviisførelsens Resultat, stolende paa, at Grundlaget
er fast. Men paa en saadan tillidsfuld Imødekommen tør, det
ligger i Sagens Natur, den Philosopherende ikke gjøre Regning.
I Philosophien kan intet Beviis staae alene; Behandlingen af en
særskilt Deel maa laane sit Lys fra det Hele, og det omfattende
Systembeviis, saa at sige, være transparent i hver enkelt
Beviisførelse. Men nu er det bekjendt af Philosophiens Historie,
at Systemerne afløse hverandre, saa at de foregaaende altid
falde for de efterfølgende. Vel er dette for en Deel ogsaa
Tilfældet i de øvrige Videnskaber, men med den rigtignok
meget betydningsfulde Forskjel, at visse Opdagelser, Methoder
og Beviismaader enten forblive uantastede eller dog under en
videre gaaende Udarbeidelse afgive Støttepunkter, kjendelige
Udgangspunkter for videre Fremskridt. Anderledes i Philosophien.
Her kan et nyt Synspunkt, en ny
Tanke, en ny Idee ikke træde ud i Verden saaledes, som en
physisk eller naturhistorisk Opdagelse. Naar et philosophisk
System med sit Princip og sin eiendommelige Methode
hjemfalder til kritisk Opløsning og taber sin Anseelse, saa
bliver det Hele saaledes omsmeltet, at man i den nye Støbning
vanskeligt kan kjende, hvad det Blivende egenlig er, som det
Forudgaaende afsatte. Betragter man Sagen udelukkende fra
denne Side, skulde man fristes til at troe, at Philosophien
ikke var nogen sand Videnskab. Men Betragtningen er eensidig
og derfor vildledende; det Vildledende ligger i, at man henter
sine Forbilleder fra andre Videnskaber og anvender en
Maalestok, som ikke passer. Philosophien maa nødvendig have
sin ideelle Eiendommelighed; hvad der væsenlig tilhører
Tankernes usynlige Verden kan umulig have samme Mærketegn,
som det, der lader sig demonstrere af stykkeviis satte
Betingelser eller paavise i synlige Kjendsgjerninger. Seet
under Tankeudviklingens og den i Klarhed voxende Bevidstheds
Synspunkt viser det sig derimod, at Philosophien i stadig,
historisk Fortsættelse arbeider paa Løsningen af visse altid
tilbagevendende Problemer, at enhver Tænker, som ved at stille
et eller flere af disse Problemer i en for Undersøgelsen
fordeelagtig Belysning har givet et Bidrag til Løsningen,
bliver staaende som en fremragende Skikkelse i Philosophiens
Historie, og at hans Indsats vedbliver at virke, om end hans
Skole for længst er ophørt. Til nærmere Forstaaelse heraf og
for, om muligt, at bidrage til en noget klarere Opfattelse af
Philosophien som Videnskab, henledes Opmærksomheden paa
efterstaaende Problemer.


% ============================================================
%  I. ERKJENDELSESPROBLEMET
% ============================================================

\section*{I.\ Erkjendelsesproblemet}
\addcontentsline{toc}{section}{I. Erkjendelsesproblemet}
% pp. 4--26  (heading printed at the foot of p.4; text runs p.4--26)
%
% Among philosophers who have treated the question of the
% conditions and limits of human knowledge, Kant occupies the
% first place. Earlier thinkers (Locke: senses; Leibniz:
% understanding) presupposed that we know from experience what
% experience is. Kant's Kritik shows that experience is not
% self-explanatory; we must ask about the conditions of its
% possibility.
%
% The pure intuitions (Rum og Tid) are a priori necessary
% presuppositions for all sensation and thereby all experience.
% The categories (including Aarsag og Virkning) are a priori
% concepts of the understanding, indispensable for all thought.
% Since pure intuitions and a priori concepts are heterogeneous
% (ueensartede), the imagination (Indbildningskraften), by
% means of temporal schematism, supplies the mediating link.
%
% The thing-in-itself (X): the experienced object arises from
% two heterogeneous contributions, one subjective (from the
% knowing subject) and one from an unknown X. The critique of
% the thing-in-itself: X, placed outside all spatial and
% temporal determination, has no bearing whatsoever on the
% filling of spatial intuition. Kant's doubling (Phænomen/
% Ting i sig) becomes a persistent equivocation.
%
% Fichte's response: dissolve the thing-in-itself; the
% Wissenschaftslehre derives the non-Ego as a necessary
% positing of the Ego's own nature-ground.
%
% Hegel's response: the Logic — concept and thing are ultimately
% the same; nature is the "being-other" of spirit, and the whole
% development is transparent in the Idea. Critique: Hegel
% presupposes what he wants to prove; the Logic presupposes Mind.
%
% Nielsen's own position: Kant is right to maintain the
% opposition between knowing subject and object (the unknown X).
% The objectifying subject cannot objectify itself; every
% objectification involves two heterogeneous factors, subjective
% and objective. The *Objektiveringslov*: as the object is, so
% its objectification must be; as the objectification is, so
% the objectifying subjectivity must be.
%
% The multiple-sensory-system argument (pp. 24--26): if there
% were beings with sensory systems A and B differing from ours,
% they would arrive at different objectifications; yet each
% system would be internally consistent. This shows the
% objectivity of experience is relational, not absolute.

% p. 4 (section heading printed at foot of p.4; text begins here)
\noindent
Blandt de Philosopher, som i nyere Tid have behandlet
Spørgsmaalet om den menneskelige Videns Betingelser og
Grændser, indtager Imm.\ Kant uimodsigelig den første Plads,
om end ikke ved de Resultater, hvortil hans Undersøgelser
% p. 5
førte ham, saa dog ved den Maade, hvorpaa han har stillet
Problemet. Forud gaaende Tænkere, hvad enten de, som Locke,
lagde Eftertrykket paa Sandserne, eller, som Leibnitz, paa
Forstanden, have uden videre forudsat, at vi af Erfaring vide,
hvad Erfaring er, at Erfaring er Noget, der forstaaes af sig
selv; men Kants „\emph{Kritik der reinen Vernunft}" viser, at
Erfaring netop ikke er Noget, der forstaaes af sig selv, at vi
for at udfinde, hvilke Erkjendelser vi kunne øse af Erfaringens
Kilder, maae gaae tilbage til Erkjendelsens ufravigelige
Grundbetingelser og begynde med Spørgsmaalet om
\emph{Erfaringens Mulighed}.

Det er indlysende, at de Forudsætninger, den Erkjendende
selv maa medbringe, for at kunne gjøre en Erfaring, ikke
kunne være laante af Erfaring, om de end først ved Erfaringens
Hjælp kunne komme til Bevidsthed. Hvilke ere da disse
Forudsætninger? Uden Sandsning er det umuligt at opfatte det
Sandselige, og uden Forstand umuligt at bestemme det Sandsede:
Erfaring er betinget af Sandsning og Tænkning; men hvoraf
er Sandsningen betinget og hvorledes forholder det sig med
Forstanden? Al Sandsning beroer paa en Anskuen. At sandse
en Gjenstand i Rummet er lige saa umuligt uden Rumsanskuelse,
som det er umuligt at opfatte en Forandring i Tiden uden
Tidsanskuelse. De rene Anskuelser: \emph{Rum og Tid}, ere altsaa
apriorisk nødvendige Forudsætninger for al Sandsning og derved
for al Erfaring. For at skaffe sig et positivt Indhold maa
Tænkningen slutte sig til Sandsningen; thi Forstanden er ingen
anskuende Evne, og den gamle Metaphysiks Forsøg paa at
aabne sig en overnaturlig Verden og ved den rene Forstand at
gjøre Opdagelser af oversandselige Ting, ere phantastiske. Men
at Sandsningen ikke kan bringe det til at gjøre nogensomhelst
Erfaring uden Forstandens Bistand, indsees deraf, at det ene
Sandseindtryk ikke kan kjendes fra det andet uden en skjel\-nende
Virksomhed, og den ene Forestilling ikke bestemmes ved
den anden uden en dømmende Virksomhed. Den dømmende
Forstand vilde ikke kunne udøve sin Virksomhed, dersom den
ikke selv \textit{a priori} indeholdt de Stambegreber, Kategorier, af
hvilke de logiske Domme nødvendig maae være betingede.
% p. 6
Hvilke disse Kategorier ere, og hvor mange de ere, kan sees
af Dommenes Inddeling i den formelle Logik. Saaledes viser
sig den hypothetiske Dom at være grundet paa Kategorierne
Aarsag og Virkning, og den disjunctive paa Kategorien
Vexelvirkning. De aprioriske Forstandskategorier ere altsaa
uundværlige for al Tænkning og derved for al Erfaring.

De rene Sandseanskuelser og de aprioriske Begreber maae
nødvendig være grundforskjellige, ja ueensartede, efterdi
Sandserne ere ueensartede med Forstanden; det Spørgsmaal
opstaaer da, hvorledes det er muligt, at saa ueensartede Former
kunne forenes i den erkjendende Bevidsthed. Vi henvises til
Indbildningskraften, som, paa den ene Side intellectuel, paa den
anden Side sandselig, ved Hjælp af Tidsanskuelsen, afgiver et
apriorisk, for Erfaringens Mulighed nødvendigt Schema. Det
er f.\ Ex.\ ikke nok at have forvisset sig om Aarsagsbegrebets
Gyldighed som Begreb alene; her fordres en Betingelse for
Begrebets Anvendelse paa Tilsyneladelserne i den sandselige
Verden, hvorved den bestemte Tidsorden faaer afgjørende
Betydning.

Ved at paavise de subjective, for al Erfaring gjældende
Grundbetingelser saae Kant sig istand til at overvinde D.\ Humes
empiriske Skepticisme. Hume forkastede Aarsagsbegrebet
som et Slags Vanebegreb, hvortil ingen Erfaring svarer. At
Tingene følge med og efter hverandre, kunne vi erfare, men at
det Ene er Aarsag til det Andet, kunne vi ikke erfare; det, vi
af Vane kalde Aarsag, er os lige saa ubekjendt som det, vi
kalde Kraft. Kant giver for saa vidt Hume Ret, som han
godtgjør det Urimelige i at ville aflede Aarsagsbegrebet af
Erfaring; men idet han med det Samme beviser, at ingen som
helst Erfaring er mulig uden Aarsagsbegreb, hævder han seier\-rig
og een Gang for alle dette Begrebs apriorisk-empiriske
Gyldighed saa eftertrykkelig, at Erfaring og Aarsag efter ham
maae siges at staae og falde med hinanden. I særlige
Erfaringstilfælde kan det ofte være tvivlsomt, hvad Aarsagen er
til en stedfindende Begivenhed. Naar A, B og C ledsages af D,
% p. 7
kan det maaskee være A, maaskee B, maaskee C, maaskee A
og B, maaskee A og C, maaskee B og C, maaskee A, B og C i
Forening, hvori Aarsagen til D maa søges; til Afgjørelse af
slige tvivlsomme Tilfælde kan man anvende Induction, men
Inductionen forudsætter netop, at der maa være en Aarsag.\footnote{%
„Zwar scheint es, als widerspreche dieses allen Bemerkungen, die man
jederzeit über den Gang unseres Verstandesgebrauchs gemacht hat, nach
welchen wir nur allererst durch die wahrgenommenen und verglichenen
übereinstimmenden Folgen vieler Begebenheiten auf vorhergehende
Erscheinungen, eine Regel zu entdecken, geleitet worden, der gemäss gewisse
Begebenheiten auf gewisse Erscheinungen jederzeit folgen, und dadurch
zuerst veranlasst worden, uns den Begriff von Ursache zu machen. Auf
solchen Fuss würde dieser Begriff bloss empirisch sein und die Regel,
die er verschafft: dass Alles, was geschieht, eine Ursache habe, würde
eben so zufällig sein, als die Erfahrung selbst; seine Allgemeinheit und
Nothwendigkeit wären alsdenn nur angedichtet und hätten keine wahre
allgemeine Gültigkeit, weil sie nicht \textit{a priori}, sondern nur auf Induction
gegründet wären\ldots\ Freilich ist die logische Klarheit dieser
Vorstellung einer, die Reihe der Begebenheiten bestimmenden Regel, als
eines Begriffs von Ursache, nur alsdenn möglich, wenn wir davon in
der Erfahrung Gebrauch gemacht haben; aber eine Rücksicht auf die\-selbe
als Bedingung der synthetischen Einheit der Erscheinungen in der
Zeit, war doch der Grund der Erfahrung selbst und ging also \textit{a priori}
vor ihr vorher." Kritik der rein.\ Vernunft. Leipzig 1838. S.~201.}

Her er, vel at mærke, ikke Tale om Forholdet imellem
Aarsagen og Tingen selv, der falder udenfor vor Bevidsthed, men
kun om Aarsagsbegrebets Nødvendighed for vor Erfaring. Et
vilkaarligt Spil af Sandsninger og Forestillinger er kun af
subjectiv Betydning og giver ingen Erfaring. Hvis en Mængde Led
a, b, c, d o.\ s.\ v.\ uophørlig forandrede deres Orden og
indbyrdes Stilling uden nogen som helst tænkelig Grund, da vilde
vi vel være nødte til at opfatte disse Forandringer i Tiden, men
Tidsfølgen vilde ikke give os noget Begreb om Aarsag og
selvfølgelig heller ingen Erfaring. Dersom Billedet af et Huus,
f.\ Ex., hvert Øieblik forandredes saaledes, at det snart vendte
Rygningen nedad, snart opad og ved Delenes Omflytning antog
Skikkelse snart af et Taarn, snart af en Hvalfisk, da maatte vi
søge Forandringernes Grund i vor egen Indbildning og kunde
paa ingen Maade henføre et saadant Billede til nogen
Erfaringsgjenstand. Der maa for Erfaringens Skyld være Noget, der
% p. 8
nøder os til at sætte det Ene forud og det Andet efter; det er,
med andre Ord, den nødvendig bestemte Tidsorden, der
udtrykker Forholdet mellem Aarsag og Virkning. Nødvendigheden
er apriorisk, thi den er en Begrebsnødvendighed; med denne
Nødvendighed bliver Gjenstanden bestemt ved Erfaringen, og
Erfaringen ved Gjenstanden.\footnote{%
„Der Begriff, der eine Nothwendigkeit der synthetischen Einheit bei sich
führt, kann nur ein reiner Verstandesbegriff sein, der nicht in der
Wahrnehmung liegt, und das ist hier der Begriff des \emph{Verhältnisses
der Ursache und Wirkung}, wovon die erstere die letztere in der
Zeit, als die Folge, und nicht als etwas, was blos in der Einbildung
vorhergehen (oder gar überall nicht wahrgenommen sein) könnte,
bestimmt. Also ist nur dadurch, dass wir die Folge der Erscheinungen,
mithin alle Veränderung dem Gesetze der Causalität unterwerfen, selbst
Erfahrung d.\ i.\ empirisches Erkenntniss von denselben möglich; mithin
sind sie selbst als Gegenstände der Erfahrung, nur nach eben dem
Gesetze möglich." Anf.\ Skr.\ S.~196--97.}

Man vil heraf forstaae, at Kant med ubøielig
Følgestrænghed maa gjøre Skilsmisse mellem Erfaringsgjenstanden
(Phænomenet) og Tingen i sig. Der maa være Noget, som
uafhængigt af os, giver vore tomme Erkjendelsesformer, de
sandselige saa vel som de intellectuelle, et virkeligt Indhold og
bestemmer vor Opfattelse af Gjenstanden; men at Gjenstanden,
saaledes som vi opfatte den, umulig kan være lig med Tingen
selv, er en ligefrem Følge af, at den ved at gaae igjennem vore
subjective Erkjendelsesmedia nødvendig maa faae et subjectivt
Anstrøg. Vor Opfattelse af Tingen er en Forandring af Tingen:
vi see Objecterne gjennem farvede Briller.

Den skarpe Adskillelse mellem Tingen i sig og dens
Tilsyneladelse finder Anerkjendelse i den exacte Videnskab, navn\-lig
i Sandseorganernes Physiologie. Helmholtz har\footnote{Handbuch
der physiologischen Optik. Leipzig 1867. S.~442--flgd.} med
Eftertryk fremhævet, at vore Billeder og Forestillinger ere
Virkninger, som de anskuede og forestillede Objecter frembringe paa
vort Nervesystem og derigjennem paa vor Bevidsthed. Her er en
Vexelvirkning mellem et Udenfra og et Indenfra, og Phænomenets
Aarsag er netop denne Vexelvirkning. At spørge, om den
Forestilling, jeg har om et Bord, dets Form, dets Fasthed, dets
% p. 9
Farve, Tyngde o.\ s.\ v., stemmer overeens med den virkelige
Ting, er det Samme som at spørge, om en Tone er rød, guul
eller blaa. Tingene bestemmes ved deres Egenskaber, og
Egenskaberne ved Vexelvirkning med vore Bevidsthedsorganer
og med andre Ting. Zinnober er rødt, vel at mærke, for dem,
der kunne adskille Farver, men sort for de Rødblinde. Uagtet
de Sidste, efter Vexelvirkningens Lov, have samme Ret i deres
Dom, som de Første, saa sige vi dog i Almindelighed: Zinno\-ber
er rødt; ligeledes tillægge vi ofte Tingen en Egenskab, som
om den tilkom Tingen i sig, skjøndt denne Egenskab kun gjæl\-der
for Tingens Vexelvirkning med andre Ting, som, naar vi
sige: Bly er opløseligt, og underforstaae: i Salpetersyre, ikke i
Svovlsyre. Resultatet er, at vore Forestillinger om Tingene kun
have praktisk Betydning. Vore Betegnelser, vore Ord og Ud\-tryk
ere kun Symboler; af disse Symboler kan der dannes et
Sprog, betegnende for vore Iagttagelser, exacte Undersøgelser
og videnskabelige Erkjendelser, men som Intet, aldeles Intet
har med Tingenes Væsen at gjøre.

Dersom dette nu er Undersøgelsens sidste, afgjørende
Resultat, dersom Erkjendelsesproblemet hermed een Gang for alle
er løst, at sige i Principet, dersom al videre gaaende Forskning
kun vil tjene til en yderligere Stadfæstelse af den kantiske
Theorie: saa maae ikke alene Fichtes, Schellings og Hegels
Arbeider betragtes som forfeilede, men ethvert fremtidigt
philosophisk Forsøg paa at belyse Tilværelsens Væsen og Grund maa
afvises som ufrugtbart og for Videnskaben vildledende.
Hvorledes det end forholder sig, saa meget er vist, at den kantiske
Erkjendelseslære har en Tvetydighed, der maa drages frem for
Lyset. Grundspørgsmaalet er: hvorledes kan det, som fylder
vore Erkjendelsesformer med \emph{reelt Indhold, være saa absolut
ueensartet med disse Former, at Tingen selv maa forblive os
aldeles ubekjendt?}

Lad os kalde det Ubekjendte X og holde os til en bestemt
Gjenstand, f.\ Ex.\ et Bord. At vi opfatte Bordet i Rummet og
tillægge det en begrændset Udstrækning, hvorved det faaer sin
Skikkelse, er indlysende; men den rene Rumsanskuelse giver
% p. 10
os intet virkeligt Bord. Bordet er ikke en Frembringelse af
vor egen Indbildningskraft; her er et Noget, der virker paa
vore Sandser, navnlig paa Synet og Følelsen, dette Noget er X.
Bordet, det virkelige Phænomen, kommer altsaa istand ved to
ueensartede Bidrag, et subjectivt fra os selv, og et objectivt fra
Tingen i sig. Vi kunne symbolisere det ved en Ligning:
$X + Y = P_{x,y}$, idet X betegner det Objective, Y det Subjective
og $P_{x,y}$ en Function af X og Y, Phænomenet. Sættes $Y = 0$
saa forsvinder det Subjective, og $X = P_x$, det i Phænomenet,
som alene skyldes det Objective, staaer tilbage. At $P_x$ ikke
selv kan være Phænomen, at der maa Subjectivitet til, for
at gjøre Tingen til Phænomen, har Kant fortræffelig viist; men
hvorledes det skeer, eller endog Muligheden af, at det kan skee,
lader han staae hen i al Tvetydighed. --- Hvorledes forholder sig
nu egenlig X til $P_{x,y}$, det af Bordet udfyldte Rum? Man kan
forestille sig Sagen paa flere Maader. X kan uden at fylde
noget Rum antages at gjøre et Indtryk paa Sjælen af den
Beskaffended, at den bestemte Rumudfyldning nøiagtig svarer til
Indtrykkets Intensitet. Det virkelige Bord er altsaa ingenlunde
et af X udfyldt Rum, thi X-Indtrykket bevirker kun, at
Subjectiviteten, der selv udfylder det givne Rum, maa bilde sig
ind, at det er en virkelig Gjenstand, der fylder det. Denne
Antagelse maa forekomme os urimelig, især naar vi betænke, at
Bordet ikke blot er en Synsgjenstand, men ogsaa en Gjenstand
for Følelsen. Overalt, hvor vi berøre Bordfladen, fornemme vi
Realiteten. Men Realiteten af enhver Sandsning hidrører jo
fra X-Indtryk; der maa da være lige saa mange forskjellige
X-Indtryk som Bordets Flade har Dele, og disse Indtryk,
hvis Forskjel her er en Stedforskjel, maae nøiagtig være
ordnede og fordeelte paa tilsvarende Maade som Fladens Dele,
d.\ v.\ s.\ ordnede og fordeelte i Rummet. Dog heller ikke denne
Antagelse tilfredsstiller, thi deraf vilde jo følge, at der til
Phænomenet af et fiirkantet Bord vilde svare, om ikke just et
fiirkantet X, saa dog X-Indtryk, der nøiagtig passe paa det
Fiirkantede, til Phænomenet af et rundt Bord X-Indtryk,
der passe paa det Runde, o.\ s.\ fremdl.\ Adskiller man de Ind\-tryk
% p. 11
af X, der gjøre Rumsbestemmelserne virkelige, fra X selv
som det, der er udenfor Rummet, da har det udenfor Rummet
værende X jo slet Intet med de virkelige Rumsphænomener at
gjøre; dets Betydning for en Udfyldning af Rumsanskuelsen er
reduceret til Nul.\footnote{%
Antager man, at X paa en vis negativ Maade reflecterer sig i sine til
Rumsanskuelsen henvisende X-Indtryk ligesom X-Indtrykkene reflectere
sig i Erfaringsgjenstandene, saa faaer man et X-Phænomen bagved ethvert
Erfaringsphænomen; men da det X, der staaer bagved X-Phænomenet,
igjen maa reflectere sig i et X-Phænomen, der endnu mere fjerner sig
fra Rumsformen, saa faaer man X-Phænomen bagved X-Phænomen i det
Uendelige, inden man naaer til det sidste, fra Rummet adskilte X, den
rene „Ding an sich".}

Lignende Vanskeligheder indstille sig, naar vi betragte de
phænomenale Forandringer, der give Tidsanskuelsen virkeligt
Indhold. Substansphænomenet viser os det i Tiden \emph{Blivende},
Accidensphænomenerne det i Tiden \emph{Vexlende}, og dog maa
Substansen charakteriseres ved Accidenserne, og Accidenserne
ved Substansen. Nu spørges: er det eet og samme X, der
bestemmer Virkeligheden saavel af Substansphænomenet som af
Accidensphænomenerne, eller er her to Slags X? Antage vi
det Første, saa indlægge vi den selvsamme charakteristiske
Forskjel i X, som i Phænomenerne. Som Substansens X
\emph{bliver} det i Tiden, som Accidensernes X \emph{vexler} det i Tiden;
dette vil da med andre Ord sige, at Tingen i sig, det
Ubekjendte, kun er istand til at fylde vor Tidsanskuelse med
virkeligt Indhold paa den ufravigelige Betingelse, at det saavel
fra Substansens som fra Accidensernes Side indgaaer i Tiden,
bliver i Tiden og forandrer sig i Tiden. Men denne Antagelse
staaer jo i aabenbar Strid med den kantiske Forudsætning.
Foretrække vi det Sidste og sætte et X for Substansen, et andet
for Accidenserne, da vil Mangfoldiggjørelsen af X ogsaa fort\-sættes
videre. Det er da ikke nok med et almindeligt X for
Substansen i Almindelighed, vi maae have et særeget X for
hver særegen Substans, et Svovl-X for Svovlsubstansen, et
Kul-X for Kulsubstansen, et Jern-X for Jernsubstansen o.\ s.\ v.;
% p. 12
hvad Accidenserne angaaer, vil det være tvivlsomt, hvorvidt vi
kunne nøies med eet X for hver Accidensgruppe, eller om der
behøves lige saa mange X'er, som der er Accidenser i hver
Gruppe. Den kantiske Fordobling: Phænomen og Ting i sig,
kommer i denne Belysning stærkt til at minde om den
platoniske: Bordet og Bordheden, Hesten og Selvhesten. Hvorledes
man end vender og dreier det, bliver følgende Spørgsmaal at
besvare: Kan X indgaae i Forandringer eller kan det ikke?
Kan det ikke, da holder det sig ganske rigtig uden for Tiden,
men saa blive ogsaa alle Forandringsphænomener til X-for\-ladte,
uvirkelige Skin. Kan X være med i Forandringerne og
give dem Præg af Erfaringsvirkelighed, da maa dette Virkelig\-hed
Givende ogsaa paa en eller anden Maade selv ingaae i
Tidsbestemmelser og være med i Tiden. At optage de X-Indtryk,
der bestemme Forandringerne, i Tidsanskuelsen, og
dog lade X være udenfor Tidsanskuelsen, er en Urimelighed,
thi det er en Modsigelse.

Den kantiske Erkjendelseslære indskærper, at vore
Forstandsbegreber kun have Gyldighed ved Bestemmelsen af
Erfaringens Gjenstande og ikke finde nogensomhelst Anvendelse
paa Tingen i sig. Men, naar Tingen i sig falder udenfor alle
Begrebsbestemmelser, hvorledes kan den da bestemmes? Den
skal ikke kunne bestemmes ved Sandsning, da den strængt
holdes udenfor Rummet og Tiden; den maa altsaa bestemmes
ved Tænkning. Da nu Forstanden umulig kan forskaffe sig
andet positivt Indhold end det, den erholder fra Sandsningen,
maa Følgen heraf jo blive, at Tingen i sig kun kan bestemmes
negativt. Man skulde nu formode, at den eneste Bestemmelse,
der turde gives af X, vilde være den, at det udelukker alle
Bestemmelser. Det maatte da være lige saa ubestemt, at det
er, som at det ikke er, lige saa ubestemt, at det gjør Indtryk
paa vore Sandser, som at det ikke gjør Indtryk o.\ s.\ v. Men i
denne Formodning blive vi høilig skuffede. X bliver udtryk\-kelig
bestemt ved en Forstandshandling; det bliver bestemt ved
en Dom, ikke ved en negativ, men ved en positiv Dom:
\emph{Tingen i sig er}. Væren, et Forstandsbegreb, bliver her,
% p. 13
tvertimod Theoriens Forskrift, anvendt paa X. Væren betegner
Realitet, noget Positivt; hvorfra faaer Forstanden dette Positive?
Ikke fra Sandsningen, thi Sandsningen giver kun det, som kan
opfattes i Rum og Tid. Heller ikke fra et Andet end
Sandsningen, thi her er jo intet Andet, der kan give Forstanden det
Positive. Saa maa X vel kjendes indirect paa sine Virkninger,
d.\ v.\ s.\ paa de Sandseindtryk, hvis Aarsag det er. Men ogsaa
denne Antagelse har Theorien udelukket, thi Aarsag er et
Forstandsbegreb lige saa vel som Virkning. Ved at bestemme
Tingen i sig som en ubekjendt, uden for Forstand og Sandsning
liggende Aarsag, en Aarsag, der skal tænkes som objectiv og
dog ikke kan tænkes objectiv, en Aarsag, der i sine Virkninger,
for Virkelighedens Skyld, maa sandses og dog ikke kan sandses,
har Kant, istedenfor at vise Erfaringens Mulighed, tvertimod
beviist Erfaringens \emph{Umulighed}.

Er der slet ingen objectiv Form, som kan tilkomme Tingen
i sig, da maa dette Ubekjendte nødvendig løse sig op i en
taaget Mulighed, en ὕλη ἄμορφος, og mangle al objectiv
Selvstændighed. Det er da saa langt fra, at den i sig ubestemte
Mulighed kan siges at bestemme vore subjective Formers
objective Indhold, at den tvertimod laaner sine tilsyneladende
Virkelighedsformer fra os og først ved dette Laan faaer Skin af
noget Objectivt. --- Her er imidlertid en Betragtning, som fra
kantisk Side kan fremhæves. Vistnok maae vi tillægge Tingen
i sig objective Former, saa sandt vi tillægge den objectiv
Gyldighed, men Sagen er, at vi slet ikke kjende disse objective
Former. Vel kan det ikke være Former som vore; men hvem
kan bevise, at der ikke gives andre Tilværelsesformer end dem,
der ligge til Grund for vor Sandsning og Tænkning? Ikke at
tale om vore fem Sandsers, Synets, Hørelsens, Lugtens,
Smagens og Følelsens charakteristiske Forskjelligheder, saa vel
indbyrdes som hos de forskjellige Individer, eller om
Muligheden af, at erkjendende Væsener kunne have baade flere og
skarpere og finere Sandser end vi, saa er der jo intet Menneske,
der kan indestaae for, at vore aprioriske Anskuelser og
Kategorier endog blot tilnærmelsesviis svare til de objective, i
% p. 14
Tingen selv udtrykte Formbestemmelser. Er ikke Rummets og
Tidens Objectivitet problematisk? Maae vi ikke, selv hvad
Erfaringsgjenstande angaaer, gjøre Forskjel mellem Tingens
Begreb og vort ufuldkomne, ofte feilfulde Begreb om Tingen?
--- Her have vi igjen Tvetydigheden.

Det falder neppe noget fornuftigt Menneske ind at nægte
Ufuldkommenheden af vor Rumsanskuelse, naar det gjælder
en klar og tydelig Opfattelse af bestemte, fleersidig combinerede
og derved indviklede Forhold i Rummet. At construere i Rum\-met
er just ikke nogen let Sag. Skulde Nogen drage dette i
Tvivl, da behøver han kun at gjøre et endog blot foreløbigt
Bekjendtskab med den analytiske Geometrie; han vil da snart
erfare, hvor let en Curve eller et mathematisk Legeme, hvis
Ligning under visse Betingelser er ham bekjendt, kan blive
ukjendelig enten ved en Forflyttelse af Coordinaternes Skjæ\-ringspunkt
eller derved, at det hele Coordinatsystem forandres.
Jo mere combineret Figuren er, desto vanskeligere ere de
specielle Relationer at gjennemskue: her kan da opstaae
mangfoldige Tvivl om Rumsforholdenes Objectivitet i det Enkelte.
Men Et er: at indrømme de af Rumsanskuelsen betingede
Forestillingers og Sandsningers Ufuldkommenhed, et Andet: at
frakjende Rumsanskuelsens aprioriske Formbestemmelser objectiv
Gyldighed. Hvis et erkjendende Væsen fra høiere Regioner
steg ned til os og erklærede, at det kunde skjelne Millioner
Dele i det Rum, vi antoge for at være opfyldt af en enkelt Deel,
da turde vi ingenlunde sige det imod; men hvis det derhos
paastod, at det tydelig kunde see Cirklen som Fiirkant og
Fiirkanten som Cirkel, mon vi saa vilde blive overbeviste om, at
dets Rumsanskuelse var mere objectiv end vor? Dersom den
Fremmede gav Prøver paa, at han ikke blot kunde sandse
Ætheren, men endog tælle 733 Billioner Svingninger i 1 Secund,
da vilde han i allerhøieste Grad forbause os; men hvis han saa
angav som Grund, at han var istand til at opfatte det
Samtidige som successivt og det Successive som samtidigt, mon vi
da deraf turde slutte, at hans Tidsanskuelse kom det Objective
nærmere end vor? Imellem Rum og Ikke-Rum, Tid og Ikke-
% p. 15
Tid ere ingen Mellemformer mulige. Hvor mange Dimensioner
Rummet kan forekomme det sandsende Væsen at have, er her
ikke Spørgsmaalet; Spørgsmaalet er om Dimension og Succession.
Ved at unddrage sig alle Bestemmelser af Rum og Tid, af
Dimension og Succession, maa hele den „objective, ubekjendte
Verden" netop blive os bekjendt som et eneste intensivt
Punkt.\footnote{At Luftbølgerne ere ueensartede med Tonerne og
Ætherbølgerne med Farverne, beviser kun, at en Aarsag i det Objective
kan blive ukjendelig ved at indgaae i Vexelvirkning med det Subjective;
men det heri Ufattelige beviser Intet hverken imod Aarsagsbegrebets eller
imod Tidens og Rummets objective Gyldighed. Om Indflydelsen af vor
Organisation skal der senere tales.}

Til Bekræftelse paa, at vore Fornuftbegreber (Ideer) lige
saa lidt finde Anvendelse paa noget i sig Objectivt som vore
Forstandsbegreber, opstiller Kant fire saa kaldte kosmologiske
Antinomier, nemlig to mathematiske og to dynamiske. Vi
holde os i denne Sammenhæng til de mathematiske. Første
Antinomie. \emph{Thesis}. Verden har en Begyndelse i Tiden, og
er, hvad Rummet angaaer, indesluttet af Grændser. \emph{Antithesis}.
Verden har ingen Begyndelse, og ingen Grændser i Rummet;
den er saavel i Henseende til Tiden som til Rummet, uendelig.
Anden Antinomie. \emph{Thesis}. Enhver sammensat
Substans bestaaer af enkelte (usammensatte) Dele, og der
existerer overalt Intet uden det Enkelte og det, som deraf er
sammensat. \emph{Antithesis}. Ingen sammensat Ting i Verden
bestaaer af usammensatte Dele, og der existerer overhovedet
intet Enkelt i den. --- Da Beviset saavel for Thesis som for
Antithesis i hver af Antinomierne maa tillægges samme
Gyldighed, saa skal denne de transcendentale Ideers indbyrdes
Strid uigjendrivelig bevise, at vi lige saa lidt kunne naae det
Objective ved Hjælp af vor Fornuft som ved Brugen af vor
Forstand. Lad os prøve Beviserne. Staaer Thesis i første
Antinomie i uopløselig Strid med Antithesis? Vi kunne ikke indsee
det. Hvis Verden er fra Evighed af, saa er der jo aldeles
ingen Modsigelse i, at der til et givet Tidspunkt kan være forløbet
en uendelig Række af successive Verdenstilstande. Vel
kan en uendelig Række ikke ved nogen successiv Synthese
% p. 16
fuldendes i en endelig Tid; den behøver en uendelig Tid, for at
fuldendes. Men fra hvilket endeligt Tidspunkt man end seer
tilbage mod den uendelig fjerne Begyndelse, seer man jo
ind i en uendelig Tid. Er der intet Modsigende i at antage
en uendelig Række af forestillede Led, saa er der heller intet
Modsigende i at antage en uendelig Række af virkelige Led,
naar Virkeligheden, vel at mærke, forudsættes at være fra
Evighed af. Hvorfor kan Verden ikke være begyndt i Tiden? Fordi,
hedder det i Antithesis, der da maatte være gaaet en Tid forud,
altsaa en tom Tid; men i en tom Tid kan Intet opstaae, efterdi der i
ingen Deel af en saadan Tid er nogen Betingelse for en Adskillen
af Tilværen fra Ikke-Væren. Men en saadan tom Tid er mere
imaginær end ellers en imaginær Størrelse. Med en imaginær Stør\-relse,
mathematisk forstaaet, kan man dog regne; men en aldeles
tom Tid, en Tid, hvori der ikke er Forskjel mellem det Forbigangne,
det Nærværende og det Tilkommende, er ingen Tid, og med
ingen Tid kan man ikke regne. Saa sandt den virkelige Ver\-den
virkelig har en Begyndelse, maa dens Begyndelse nødvendig
falde i en virkelig Tid, og den virkelige Tid begynde med
den virkelige Verden. Da det i denne Henseende ingen
Forskjel gjør, om vi rykke Begyndelsens Tidspunkt en halv eller
en heel Evighed tilbage, saa er der jo, hvad Verdens Begyndelse
angaaer, ingen Strid mellem Thesis og Antithesis. --- Ved
at betragte Rummets Begrændsning og Ikke-Begrændsning
komme vi til samme Resultat. Den virkelige Verden er begrændset,
thi det udvortes Virkelige maa i ethvert Øieblik
have bestemte Grændser og for saa vidt være et Endeligt. Men
Verden i dens Endelighed og Begrændsning befinder sig ogsaa
i virkelig Vorden, og kan saaledes i hvert Øieblik være i Begreb
med at overskride enhver endelig Grændse. Det er den
virkelige Verden, der gjør Rummet objectivt og sætter objective
Grændser i Rummet; Verdensrummet har ingen særegen Væren
udenfor Verden. Antithetiken er bygget paa Umuligheden af
at fuldende den successive Synthese af Leddene i en uendelig
Række, d.\ e.\ paa Umuligheden af at sammenfatte alle Dele af
en uendelig Totalitet. Men, spørge vi, vil da et begrændset
% p. 17
Hele, der kun indeholder 7 Dele, være mere virkeligt end et
Hele, der indeholder 7 Billioner Dele, fordi den successive
Synthese i første Tilfælde falder os let, men i sidste Tilfælde
umulig? Det vil dog vel Ingen paastaae. Men, hvis vi tør
antage en Totalitet med 7 Billioner Dele, et Verdenshele med
7 Billioner Firmamenter, for virkelig, uagtet den successive
Synthese kræver en uhyre lang Tid, hvor kan der saa være
nogen Modsigelse i at forudsætte Virkeligheden af en uendelig
Totalitet, fordi en successiv Synthese (rettere: Tilblivelse) af
dens enkelte Dele maae have udkrævet en uendelig lang Tid?

For \emph{Thesis} i den anden Antinomie anføres, at der intet
Substantielt kunde gives, dersom der ingen udeelbare Dele
existerede; det substantielt Sammensatte maa altsaa bestaae af
usammensatte Dele. Til Begrundelse af \emph{Antithesis} anføres, at
ingen Sammensætning af Substanser er mulig uden i Rummet.
Den i Rummet existerende Substans maa nødvendig have
substantielle Dele; men Rummet har ingen substantielle Dele, det
er et \emph{Totum}, intet reelt \emph{Compositum}. Modsigelsen skal nu
være: at sammensatte Substanser ikke kunne existere uden
Enkeltsubstanser, og at Enkeltsubstanser ikke kunne existere i
Rummet. Istedenfor at bestaae af Enkeltsubstanser kommer
den sammensatte Substans til at bestaae af sammensatte
Substanser i det Uendelige, og det er urimeligt. --- Denne
Modsigelse er imidlertid mere imaginær end virkelig. Ere Enkelt\-substanser
og Rum uforenelige Størrelser, saa er ikke blot
Erfaringen, men ogsaa Mathematiken umulig, thi saa er
Anskuelsen og dermed Sandsningen absolut uforenelig med Tænk\-ningen.
„Det er her", siger Kant, „ikke nok til det rene
Forstandsbegreb af det Sammensatte at finde Begrebet af det
Enkelte, man maa finde den Anskuelse af det Enkelte, der
svarer til Anskuelsen af det Sammensatte, og dette er, efter
Lovene for Sandseligheden, altsaa ogsaa ved Sandsegjenstande,
aldeles umuligt." --- Dette Umulige er imidlertid ikke saa
indlysende, som det ved første Øiekast kunde synes. Et
mathematisk Legeme, en Kubus f.\ Ex., er en Gjenstand for
Anskuelse; hvorledes forholder det sig her med det Enkelte og det
% p. 18
Sammensatte? Den givne Kubus kan anskues som et
Enkeltværende, et af 6 qvadratiske Flader begrændset Continuum;
men den kan ogsaa betragtes som et Sammensat og opløses i
Elementer. Betegnes den ved en uendelig Summation af
forestillede Legemlighedselementer $x^2dx$, saa er Trediedelen af den
anskuede Kubus et ved endelige Grændser bestemt Integral af
disse Elementer. Fladen kan opløses i Fladeelementer $xdx$,
Linien i Linieelementer, $dx$, og disse igjen i Elementer af
høiere Orden o.\ s.\ fremdl.\ I Forhold til Legemet er
Legemlighedselementet det Enkelte, ligesom Fladeelementet er det
Enkelte i Forhold til Fladen, og Linieelementet det Enkelte i
Forhold til Linien. Tanke og Anskuelse ere, hvad det
mathematiske Legeme angaaer, enige om, at det Enkelte maa være
relativt for det Sammensattes Skyld, og det Sammensatte
relativt for det Enkeltes Skyld, at altsaa Fastsættelsen af et absolut
for sig værende Enkelt ikke er en Fornuftidee, men en Fixidee.
Mon det Samme ikke skulde gjælde for det materielle Legeme?
Det skal, ifølge Kant, være en Urimelighed, „foruden det
mathematiske Punkt, der er enkelt, men ingen Deel, kun en Grændse
i Rummet, at tænke sig physiske Punkter, der vel ere enkelte,
men dog have det Fortrin, som Dele af Rummet at udfylde
dette ved deres Aggregation." Men den foregivne Urimelighed
vil ikke blive os indlysende. Lad os kun sætte Atomer og
Moleculer istedenfor mathematiske Elementer, og Aggregationer
istedenfor uendelige Summationer: Grundforholdet mellem
det Enkelte og det Sammensatte vil dog blive det samme. At
den endelige Discretion er en Udstykning, et Brud paa
Continuiteten, maa indrømmes. Men om end Atomerne ere, i reel
Forstand, nok saa udeelbare, ja selv, om de, hvad der af andre
Grunde er en Urimelighed, antages at være absolute Enkelt\-substanser,
saa kan dog Rumsanskuelsen ikke have det mindste
herimod at indvende. Atomets substantielle Enkelthed staaer
ingenlunde i Strid med dets mathematiske Deelbarhed. Den i
Rummet udstrakte Enkeltsubstans har ganske vist lige saa
mange Substanselementer, som det Rum, den fylder, har Rums\-elementer,
d.\ v.\ s.\ den har uendelig mange, thi Discretionen
% p. 19
gaaer i det Uendelige. Men da den uendelige Discretion, hvad
udtrykkelig sees af Rummet, jo falder sammen med Continuiteten,
hvad kan saa Anskuelsen indvende imod, at Enkeltsubstansen,
et reelt Continuum, udfylder Rummet, det formelle Continuum?
Idet Udstykningen er forudsat, træder Aggregationens
Endelighed istedenfor Summationens Uendelighed; men da den
virkelige Afstand mellem virkelige Dele netop giver Rumsformen
Præg af Virkelighed, maa Aggregationen jo være et talende
Beviis for Rummets formelle Objectivitet.

Kant er fuldt forvisset om sine Antinomiers Uopløselighed,
„denn für die Richtigkeit aller dieser Beweise verbürge ich
mich\footnote{Prolegomena. Werk.\ 3ter Band. S.~264.};" men
Uopløseligheden forudsat, hvad er saa herved vundet? „At vi ikke
maae tænke os Sandseverdenens Phænomener som Ting i sig
selv." Det maa indrømmes. Men der følger endnu noget Andet
deraf, Noget, som det er betænkeligt at indrømme; thi naar det
udtrykkelig tilføies, at vi heller ikke maae tænke os
Grundsætningerne for Phænomenernes Forbindelser som gjældende for
Tingene i sig, saa gaaer Forstanden tabt, saa sortner det for
Fornuftens Øie. Dersom vi endda kunde blive staaende ved det
foreløbige Resultat, at vi i al vor Erkjenden maae holde os til
vore Erfaringer og lade det uafgjort, hvorledes Tingene selv i
den os ubekjendte Verden egenlig ere beskafne: det fik at gaae.
Men Antinomierne drive os til en Yderlighed, som enten maa
slaae over i bundløs Skepticisme eller vække Fortvivlelse. Vi
skulle antage Ting i sig, som, uagtet de betinge vore Erfaringer,
dog udelukke alle Erfaringsbetingelser, idet hverken Rum eller
Tid, hverken Eenhed eller Fleerhed, hverken Grund eller Følge,
hverken Aarsag eller Love, hverken Heelhed eller Dele, hverken
Enkelthed eller Sammensætning have mindste Betydning for
disse i sig usandselige og utænkelige Ting. Det er haarde Vilkaar.

Til Forstaaelse af Fornuftideernes Strid i Antinomierne
anfører Kant Følgende: „To hinanden modsigende Sætninger
kunne ikke begge være falske, med mindre det Begreb, der lig-
% p. 20
ger til Grund for dem begge, i sig selv er modsigende; f.\ Ex.
de to Sætninger: en fiirkantet Cirkel er rund, og en fiirkantet
Cirkel er ikke rund, ere begge falske." Men naar den
menneskelige Bevidsthed kun har Valg mellem Et af To: enten at
lade Fornuftideerne eller „Tingene i sig" være „fiirkantede
Cirkler", saa er det i det Mindste undskyldeligt, om den endnu
engang fordrister sig til at hævde Fornuftens Ret ved at vælge det
Sidste. Dette Valg er charakteristisk for J.\ G.\ Fichtes
„Wissenschaftslehre".

Vi kunne i al Korthed betegne Afvigelsen fra det kantiske
Synspunkt ved at ombytte Symbolet $X + Y = P_{x,y}$ med
$XY = P$. At dette nødvendig maa give Undersøgelsen en anden
Retning og lede til et heelt andet Resultat, er indlysende.
Naar Y, det Subjective, bliver Nul, d.\ v.\ s.\ naar den
erkjendende Subjectivitet, Jeg, forsvinder, bliver det hele til Nul:
intet ubestemmeligt Object, intet for sig existerende X, staaer
tilbage. Ved at fastsætte Vexelbestemmelsen af Subjectivt og
Objectivt, Jeg og Ikke-Jeg, har Fichte indført en stor
Grundtanke i Philosophiens Historie, en Fundamentalsætning af saa
indgribende Betydning, at den lige saa lidt lader sig udelukke
af Erkjendelseslæren som Tyngdeloven af Physiken. Endskjøndt
Jeg og Ikke-Jeg gjensidig udelukke hinanden, forstaaer det
sig dog af sig selv, at X, Ikke-Jeg, umulig kan tillægges nogen
udenfor Jeget gjældende Selvstændighed, saa lidt som det kan
smelte sammen med Jeget selv. Som real Modsætning til Jeget
maa Ikke-Jeget altsaa være et Element i Jegets Væsen, en
Skranke, der er uundværlig for Jegets Selvudvikling. Kun
derved, at det tænkende Jeg bliver sig sin Modsætning, sin
Skranke, bevidst som en ubevidst Forudsætning, der er
fremgaaet af dets eget Væsen, bliver det sig selv bevidst, og
overvinder Skranken theoretisk. Erfaringsgjenstandene, de udvortes
Phænomener, fremstille sig umiddelbart og med en Nødvendighed,
hvorover Jeget ikke raader, og antage derved et Præg
af udvortes Realitet; men at Jeget ikke alene erkjender de
Former og Kategorier, hvori Gjenstandene fremstille sig, for
sine egne Former, men tillige erkjender Nødvendigheden, hvor-
% p. 21
med de fremstille sig, for en Nødvendighed af Jegets egen
Naturgrund, er i „Videnskabslæren" Erkjendelsesproblemets Løsning.

Det kritiske Blændværk hos Kant bestaaer i, at den
ubekjendte Ting bestandig tager sig ud som noget i sig Selvstændigt,
Objectivt, Positivt, ja som det eneste ubetinget Positive,
uagtet det fra alle Sider viser sig, at dette Positive kun kan
bestemmes negativt. Det Negative er positivt, det Positive er
negativt; denne Modsigelse, mod hvilken Kritikeren af al Magt
værger sig, dukker allevegne op og driller Kritiken. Dette
Blændværk har Fichte med klar Bevidsthed gjennemskuet, idet
han har indseet det Negatives Betydning og tillagt X en
væsenlig Værdi som Ikke-Jeg.

Men, trods dette umiskjendelige Fremskridt er den
fichtiske Løsning lige saa utilfredsstillende som den kantiske.
Hvad der fra objectiv Side er umuligt for Kant, nemlig at
paavise, hvorledes det i sig ubestemmelige X kan fylde
Erkjendelsesformerne med en uoverskuelig Mangfoldighed af reale
Bestemmelser, det Samme er fra subjectiv Side umuligt for
Fichte. Som Ikke-Jeg er X altfor eensformigt til at afgive
nogen for Phænomenernes Forskjellighed tilstrækkelig
Forklaringsgrund; en Steen er et Ikke-Jeg, en Plante et Ikke-Jeg, et
Dyr ligesaa: Dyr og Planter have i Forhold til Jeget samme
Betydning som Mineraler og Stene. Den subjective Reflexion
af Jeg og Ikke-Jeg fortærer alle objective Relationer mellem
Phænomenerne indbyrdes. Kant faaer dog en Erfaringsverden,
om end paa et tvetydigt, sig selv opløsende Grundlag. Fichte,
der ophæver Grundlagets Tvetydighed, faaer ikke engang
Saameget; han faaer kun en subjectiv nødvendig Skyggeverden af
Skranker og Ikke-Jeg, en Verden saa tom, at Naturindholdet
forsvinder, og den philosophiske Interesse for Erfaringens
Realiteter gaaer tabt.

Den hegelske Philosophie lader Erkjendelsesproblemet gaae
under i den absolute Viden istedenfor at blive løst.
Modsætningen af den erkjendende Subjectivitet og Erkjendelsens
Object, der er Problemets Bærer, løser sig op i den dialektiske
Eenhed af Tænken og Væren, der gjennemstrømmer Logiken
% p. 22
og udmunder i Ideen. I Ideen, det Absolute, er Subjectiviteten
ophævet, d.\ v.\ s.\ den er forvandlet til Sidebestemmelse, Moment,
om end et Totalitetsmoment; og det Samme gjælder om
Objectiviteten, dog med kjendelig Overvægt over det Subjective.
Istedenfor et anskuende, tænkende Selvvæsen, der har Begrebet
til Indhold og Begrebets Realitet til Gjenstand, faae vi hos
Hegel et Begreb, der begriber sig selv, en Idee, som efter at have
givet sig selv Subjectivitet i Reflexionens Form, igjen ophæver
denne sin Subjectivitet og omsætter sig i Umiddelbarhedens
Form som Objectivitet i tilværende Objecter. Begrebet, i
absolut Eenhed med sit Indhold \emph{Ideen}, kan umulig holde sig i
to modsatte Former paa een Gang; naar det Absolute er sat
som subjectiv Idee i Subjectivitetens d.\ e.\ i Reflexionens Form,
er det Objective nedsat til Moment og forsvundet; naar det
Absolute er sat som objectiv Idee i Objectivitetens d.\ e.\ i
Umiddelbarhedens Form, er Subjectiviteten forsvunden: det
maa da begribe sig selv. Men naar det Objective skal kunne
begribe sig selv lige saa vel som det Subjective, saa er det sig
begribende Begreb ubegribeligt. Skal den absolute Idee være
Slutningseenhed af subjectiv og objectiv Idee, da maa den enten
antage en Form, som, idet den hverken er subjectiv eller
objectiv, hæver den over Extremerne, eller det maa være
en af hine to modsatte Former, hvori den holder Extremerne
sammensluttede. Det Første er, fra Hegels Standpunkt,
uantageligt; thi Schellings Indifferens af Subjectivt og Objectivt
er kun en abstract Forudsætning, en Ugrund, der staaer langt
under Ideerne, og et Oversubjectivt eller et Overobjectivt har
ingen Hjemmel i „den absolute Viden". Med Hensyn paa det
Sidste maae vi da vælge enten det Subjective eller det Objective
til oprindelig Form for den absolute Eenhed; men den hegelske
Philosophie gjør os Valget umuligt. Sammenholder man Logiken
med Naturphilosophien, saa tager det sig ud, som var
Naturideen, Ideen i Umiddelbarhedens objective Form, det Afledede,
og Logikens Idee, Ideen i den subjectivt reflecterede Form, det
Oprindelige. Men ved at sammenholde Naturphilosophien med
Aandsphilosophien overbevises vi om, at Sagen forholder sig
% p. 23
omvendt; thi Logik i logisk Reflexion forudsætter Aand. Der
kan umulig være logisk Reflexion, hvor der ingen Aand er;
men Naturen er Aandens Forudsætning, ligesom Umiddelbar\-heden
er Reflexionens Forudsætning. Naturen er Aandens
„An sich", og Naturen er det Objective; det er altsaa det
objective Begreb, der --- paa ubegribelig Maade! --- begriber det
Subjective. At Hegel ved de ideelle Formbestemmelsers
immanente Dialektik gjør et Kjæmpeskridt fremad, at han ved at
gaae videre end Schelling faaer Negativiteten udviklet efter en
anden Maalestok og i en ganske anden Dybde end enten Kant
eller Fichte, maa indrømmes; men det Absolute selv faaer han
ikke klaret. Trods al Talen om den sig tænkende Tænken,
den sig absolut vidende Aand, er det umuligt at faae at vide,
hvad det Vidende selv, det absolut Vidende, er, om det er
absolut subjectivt eller absolut objectivt, om det er et Selv eller
en selvløs Proces eller maaske ingen af Delene. Det er et sig
begribende Begreb, der skal klare Alt og bringe den sidste
Rest af det Ubekjendte til at forsvinde; men langt fra at have
elimineret X, har det sig ubegribeligt Begribende forvandlet
det Absolute selv til et X, en uigjennemtrængelig Gaade, der
passende kan betegnes ved et gaadefuldt Symbol: $XY = X_{x,y}$.

Skal Erkjendelsesproblemet kunne løses, maae følgende
Grundbetingelser ufravigelig fastholdes:

\begin{enumerate}
\item[1)] \emph{Der maa være en uopløselig Modsætning af
  Subjectivt og Objectivt, hvori Subjectiviteten selv med sine
  aprioriske Former er overgribende Eenhed.}
\item[2)] \emph{Erkjendelsen maa beroe paa en saa
  gjennemgaaende Vexelbestemmelse af det Subjective og det Objective,
  at intet Object kan enten være eller tænkes at være uden
  tilsvarende Objectivering.}
\item[3)] \emph{Den Objectiveringens Ufuldkommenhed, der har
  sin Grund i de existerende Objecters Uafhængighed af organisk
  betingede Subjecter, maa lade sig afhjælpe ved subjective
  Æqvivalenter af ontologisk Objectivering.}
\end{enumerate}

% p. 24
Imod Fichte og Hegel har Kant, hvad den første Grundtanke
angaaer, umiskjendelig Ret, naar han fastholder Modsætningen
mellem det erkjendende Subject og Erkjendelsens Object,
ligesom han da ogsaa --- fra Endelighedens Side --- har Ret i
at forudsætte et Ubekjendt, en Ting i sig, som ikke gaaer op i
vor Erkjendelse. Men i den Maade, hvorpaa han bestemmer
Tingen i sig, har han høiligen Uret; thi det er ikke ved at
udelukke de apriorisk subjective Formbestemmelser, at Tingen
unddrager sig Erkjendelighed\footnote{Den nyere Kritik er imidlertid
gaaet rask frem, idet den har udstrakt Indflydelsen af „vor Organisation"
til Kategorier som Eenhed og Fleerhed, Substans og Accidens, Subject og
Prædicat o.\ s.\ v. Vi ere, hedder det, naturligviis nødsagede til at
anvende slige Kategorier paa Grund af vor Forstands eiendommelige
Organisation; men deraf følger endnu ikke, at de i sig selv ere
uforanderlige. Gaadefuld Kritik! At Forstanden ikke tør opfattes som en
Evne for sig over for de andre Sjæleevner, at dens Virksomhed maa være
afhængig af organiske Dispositioner, kan sees af forskjellig Begavelse;
Individets Forstand kan være oplagt til en Art intellectuel Virksomhed
uden i samme Grad at være oplagt til en anden. Heller ikke kan det
nægtes, at en eensidig Forstandsretning kan forlede til at fastsætte
Grundbegreber som exclusivt aprioriske, der dog vitterlig ikke ere det;
man tænke paa Pythagoræernes philosophiske Opfattelse af det dekadiske
System f.\ Ex. Men slige Feilgreb blive jo at corrigere af Forstanden selv
efter Forstandens Love. De Kritikere derimod, der give sig Mine af at
kunne tænke sig --- ikke blot en Forstand, der er uendelig høiere end den
menneskelige, thi det kunne vi alle --- nei, en Forstand, der annullerer
Forskjellen mellem Eenhed og Mangfoldighed, Aarsag og Virkning, Subject
og Prædicat, maae tilvisse udmærke sig ved en ganske eiendommelig
Organisation. „Metaphysiken" kunde de nu sagtens afskaffe; men hvordan
gaaer det med Forstanden? For at tænke sig en Forstand, der er absolut
ueensartet med den menneskelige Forstand, maa Mennesket jo gaae fra
Forstanden.}, men derved, at dens Opfattelse
modificeres efter Sandsernes og de sandsende Væseners
subjective Eiendommelighed. Den ubekjendte Ting, der er Gjenstand
for Sandsning, er en Ting med sandselige Egenskaber lige saa
vel som det Erfaringsphænomen, der ad Sandsningens Vei
kommer til vor Bevidsthed. I Sandsesystemer af forskjellig
Orden, A, B, C o.\ s.\ v., maa det samme ubekjendte X med sine
Egenskaber nødvendig fremstille sig i forskjellige Phænomener.
I A-Systemet ville Egenskaberne $x_1, x_2, x_3, \ldots$ præsentere sig
% p. 25
som $a_1, a_2, a_3\ldots$, i B-Systemet som $b_1, b_2, b_3\ldots$ og
saaledes i det første Tilfælde give Erfaringsgjenstanden $P_a$,
i det sidste Tilfælde $P_b$. De erkjendende Væsener have i det
ene System lige saa megen Ret, som de i det andet. Men
naar de erkjendende Væsener alle maae indrømme, at Tingen
selv falder uden for deres Erkjendelse, hvorpaa beroer da
Erkjendelsens Sandhed? Besvarelsen af dette Spørgsmaal fører
os til den anden Grundbetingelse.

Dersom Objectet i sig med sine Egenskaber, sine concrete
Bestemmelser virkelig kunde staae selvstændig udenfor al
Objectivering, og dersom det saaledes staaende Objective var det
virkelig Sande, saa vilde Sandhedens Erkjendelse være umulig.
Men ved at gjennemarbeide Forholdet mellem Jeg og Ikke-Jeg
har Fichte slaaende beviist, at det Objective, langt fra at være
selvstændigt, er afhængigt af og bestemt ved Subjectiviteten.
Den Overdrivelse, hvormed han hævder det rene Jegs
oprindelige Selvindskrænkning og den Energie, hvormed det
selv sætter sine Objecter, afgiver en høist fornøden Modvægt
mod den vulgære Antagelse, at Objectiveringen egenlig ikke
vedkommer Objecterne, men at det bliver Subjectivitetens egen
Sag, hvorvidt den objectiverer sandt eller falsk. Naar vi nu
indrømme Fichte, at ethvert virkeligt Object maa være et
nødvendigt Produkt af en objectiverende Subjectivitet, men umulig
kunne indrømme ham, at Objecterne virkelig i sig selv ere
saaledes, som vi objectivere os dem, eller at vor menneskelige
Objectivering skulde formaae at give Tingene Existens, saa tye vi
til den tredie Grundbetingelse, for at udfinde, hvorledes den
endelige Objectiverings Ufuldkommenhed bliver at afhjælpe.

Erkjendelsens almeengyldige Grundlag falder sammen med
Objectiveringens ontologiske Grundlag. Naar En f.\ Ex.\ foredrager
Hypothesen om den roterende Urdamp, hvoraf vort Planetsystem
skal have udviklet sig, da antager han visselig ikke, at hiin
Damp har været et Erfaringsobject for menneskelig Sandsning,
saa lidt som hans Mening er, at hvad der den Gang var „en
Ting i sig", siden hen ved vore Erfaringer er blevet forvandlet
til Damp. Heller ikke antager han, enten, at Urdampen
% p. 26
existerede i vor subjective Rumsanskuelse og roterede i vor
subjective Tidsanskuelse, ei heller at det er os, der med
apriorisk Nødvendighed have foreskrevet Naturen de Love, hvorefter
den hele Udvikling er foregaaet. Nei han antager vel snarere,
at det altsammen, Rummet, Tiden, Dampen, Bevægelsen, har
været objectivt uden nogensomhelst Objectivering. Da det nu
er beviist, at den sidste Antagelse er lige saa meningsløs som
den første, hvad følger saa heraf? At her i Stilhed maa være
underskudt et subjectivt Æqvivalent af ontologisk Objectivering!
Hypothetisk bliver dette udtrykt saaledes: dersom vi havde
været tilstede og kunnet overskue det Hele, da maatte Sagen
nødvendig have taget sig ud, som Beskrivelsen lyder. Altsaa: hver
Gang, vor indskrænkede Subjectivitet gaaer ud over sine egne
Erfaringsgrændser, bliver den for det Objectives Skyld nødsaget
til at forudsætte en Anskuen, som ikke er vor Anskuen, en
Seen, som ikke er vor Seen, en Tænken, som ikke er vor Tænken,
en Objectivering, som ikke er vor Objectivering; den bliver,
kort sagt, nødsaget til at forudsætte en ontologisk objectiverende
Subjectivitet.

Af Objectiveringens ontologiske Grundlov, symboliseret
ved $XY = O$, idet O betegner Objectet, følger: at ethvert
objectivt Moment nødvendig maa forudsætte et tilsvarende
subjectivt, efterdi $Y = 0$ giver $O = 0$, og at omvendt ethvert
subjectivt Moment i en virkelig Objectivering maa være bestemt
af et tilsvarende objectivt, efterdi $X = 0$ giver $O = 0$.
Subjectivt og Objectivt ere, i ontologisk Forstand, correlate
Begreber, men ikke coordinerede: den objectiverende Subjectivitet er,
som Modsætningens overgribende Eenhed, et \textit{dignitate prius}.

Af Grundloven følger endvidere, at Erkjendelsens Sandhed
hverken beroer paa det Objective for sig ei heller paa det
Subjective for sig, men paa Vexelbestemmelsen af begge. Lad der
end være Tusinder, ja Millioner af forskjelligartede
Sandsesystemer i Universet, saa ere Erkjendelsens aprioriske Former
dog den almene Fornufts subjective Grundformer, og den almene
Fornuft er fælles for alle erkjendende Væsener. De erkjendende
Væsener, der med samme Sandsesystem høre ind under samme
% p. 27
Orden, maae indsee, at den concrete, af Sandsningen betingede
Objectivering, har al den Sandhed, den kan og skal have, naar
Vexelvirkningen mellem det Subjective og det Objective er for
vedkommende Orden nøiagtig bestemt. De uligeartede
Sandsesystemer ere lige saa lidt istand til at vække Tvivl om det
Almeengyldige i Objectiveringens Eiendommelighed, som den
Rødblinde, trods Sandheden af hans Objectivering, er istand til at
vække Tvivl om Farvernes Forskjellighed hos dem, der have
Farvesands. Sandheden er, at Vexelvirkningen mellem Subjectivt
og Objectivt maa være forskjellig for Subjectiviteter af
forskjellig Orden; men \emph{Vexelvirkningen selv er Sandheden}. At
antage et udenfor Vexelvirkningen staaende, uobjectiveret X for
det virkelige Sande, er en Drøm, en Indbilding; udenfor
Vexelvirkningen er der ingen Sandhed, thi udenfor
Vexelvirkningen er der, ontologisk forstaaet, slet Intet.


% ============================================================
%  II. REALITETSPROBLEMET
% ============================================================

\section*{II.\ Realitetsproblemet}
\addcontentsline{toc}{section}{II. Realitetsproblemet}
% pp. 27--50
%
% The problem only becomes properly formulated in the explicit
% opposition between Positivism (Comte) and absolute Idealism
% (Hegel). Comte: there is no Absolute, only existing things
% in relation. Hegel: the only thing that truly is, is the
% Absolute; all relations have their subsistence in an
% encompassing unity. Both have their element of truth.
%
% The problem: Thing and concept must be assumed to be both
% the same and yet not the same.
%
% That they are the same: any thing's full set of
% determinations just is its concept. Chemistry: the specific
% affinities and molecular structure of gold are expressions of
% its real concept; the physicist's grasp of laws is not
% limited to external relations but penetrates into each
% thing's inherent system of law-determinations.
%
% That they are not the same: the immediate existence of the
% thing cannot be reduced to its logical form; things are real,
% immediate, "anti-logical," their relations are power-relations.
%
% The Objectiveringslov applied to reality: the positivists'
% error is to rush directly at facts and objects without
% examining the relation between object and objectification.
% The idealists' error is to dissolve the objective into the
% logical. The correct position: objectification of power-objects
% requires a power-objectification; the real concept of a thing
% is its law-determined power-structure, not a logical abstraction.
%
% Teleology vs. causality: if only physical causes are admitted,
% final causes can have no objective significance. But the
% Objectiveringslov shows that teleological descriptions cannot
% be eliminated even from biology, since the living being has
% a "Selvhedsbegreb" (a self-concept with surplus of ideality)
% irreducible to a "Stofbegreb" (a concept of matter).
% Section II closes (p.50) with the transition to the two
% relatively independent causalities, ideal and real — the
% "two heterogeneous principles" that name section III.

% [text to be added]


% ============================================================
%  III. DE UEENSARTEDE PRINCIPERS PROBLEM
% ============================================================

\section*{III.\ De ueensartede Principers Problem}
\addcontentsline{toc}{section}{III. De ueensartede Principers Problem}
% pp. 50--77   (heading printed on p.50, PDF 59; runs to the end, p.77)
%
% [Heading VERIFIED against the scan: "III." / "De ueensartede
% Principers Problem." — this single section subsumes the
% soul--body problem, the problem of freedom, and Tro og Viden.]
%
% The soul--body problem (opens p.50): how the Subjective in
% reality relates to the Objective; the living organism as a
% "betinget Selvvirksomhed" (conditioned self-activity); the
% distinction between the lifeless and the living; the organic
% unity as a form of ideal causality under material conditions;
% Stofskifte og Formgivning; the ideal midpoint that cannot be
% mathematically localised nor reduced to physico-chemical
% summation.
%
% The problem of freedom: freedom is not unconditioned, but
% "selvbevidst, betinget Selvvirken" (self-conscious, conditioned
% self-activity). Motiv og Villie: from one side motives
% determine the will as real causes; from the other the
% self-conscious will raises itself above the web of motives
% by its own act. Neither determinism nor indifferentism is
% adequate. The ideal--real opposition in ethics.
%
% Tro og Viden (from ~p.70): the choice between opposed
% possibilities (Religion og Videnskab, Tro og Viden). The key
% statement (p.\ 72): "Principerne for Tro og Viden ere absolut
% ueensartede." Not a counsel to sever religion from science or
% faith from reason: the dualistic separation is itself the
% condition for a reconciling unity. Miracle-critique
% (Troeskritik vs. Videnskritik); Aandsvirkelighed. Closing
% statement (p.\ 77) — verify verbatim against the scan.

% [text to be added]

\end{document}
